%% sim_methods.tex
%% Methods section: simulation description
%% Covers GIZMO/MFM, FIRE-3 cosmological setup, and STARFORGE Jeans-refinement.
%% Intended to be \input{} into the main paper or compiled standalone.
%%
%% Key citations needed in the main .bib file:
%%   Hopkins2015   – GIZMO methods paper  (MNRAS 450, 53)
%%   Hopkins2018   – FIRE-2 methods paper  (MNRAS 480, 800)
%%   Hopkins2022   – FIRE-3 methods paper  (arXiv:2203.00040)
%%   Grudic2021    – STARFORGE I methods   (MNRAS 506, 2199)
%%   Grudic2022    – STARFORGE II          (MNRAS 512, 216)
%%   Wetzel2023    – FIRE public data release
%%   Springel2005  – tree-PM gravity (Gadget-2)
%%   Planck2015    – cosmological parameters

\section{Simulation}
\label{sec:simulation}

\subsection{Numerical method}
\label{sec:gizmo}

All simulations are performed with the \textsc{Gizmo} code
\citep{Hopkins2015}, using the Meshless Finite Mass (MFM) method for
magnetohydrodynamics.  MFM is a quasi-Lagrangian, Galilean-invariant
finite-volume scheme in which the fluid equations are solved on a set of
moving mesh-generating points whose masses are conserved exactly (no
inter-particle mass flux), while fluxes of momentum, energy, and magnetic
flux are exchanged through the effective faces of a Voronoi-like
tessellation reconstructed at each timestep from the particle distribution.
Spatial gradients are estimated via a least-squares matrix inversion over the
kernel-weighted neighbour list ($N_\mathrm{ngb} = 32$), and a second-order
slope-limited reconstruction is applied before each Riemann solve.  This
combination gives \textsc{Gizmo}/MFM second-order convergence in smooth
flows, near-zero numerical diffusion in shear flows, and automatic
adaptivity of the resolution element to the local fluid density -- key
advantages over traditional SPH for problems involving shocks, turbulence,
and gravitational collapse.

Gravity is solved with a hierarchical oct-tree algorithm combined with a
particle-mesh (PM) long-range force, following the approach of
\citet{Springel2005}.  Adaptive gravitational softening is used for gas
particles, with the softening length tied to the kernel radius so that
force and hydro resolution are matched throughout the simulation.  Sinks
(PartType5) use a fixed gravitational softening equal to the sink accretion
radius.  All particles are advanced with individual, adaptive timesteps
using a Courant--Friedrichs--Lewy (CFL) criterion for hydrodynamics and a
gravitational criterion based on the local tidal tensor \citep{Grudic2021}.

The ideal-MHD equations are evolved together with the fluid equations using
a constrained-gradient divergence-control scheme
(Powell--Dedner cleaning) to maintain $\nabla\cdot\mathbf{B}\approx 0$ to
machine precision.  Non-ideal MHD terms (ambipolar diffusion, Ohmic
resistivity, the Hall effect) are available in \textsc{Gizmo} via
operator-splitting and are relevant at sub-AU scales beyond the resolution
of the present simulation.

\subsection{Cosmological simulation: the \textit{m12f} zoom-in}
\label{sec:m12f}

The host simulation is the \textit{m12f} cosmological zoom-in from the FIRE
suite \citep{Hopkins2022,Wetzel2023}.  The zoom-in targets a Milky
Way-mass dark matter halo ($M_{200}\sim 10^{12}\,\Msun$ at $z=0$) drawn
from a periodic box of comoving side length $\sim\!85\,\mathrm{Mpc}$, with
initial conditions based on the Planck-2015 cosmological parameters
($\Omega_\mathrm{m}=0.3089$, $\Omega_\Lambda=0.6911$,
$\Omega_\mathrm{b}=0.0486$, $H_0 = 67.74\,\mathrm{km\,s^{-1}\,Mpc^{-1}}$,
$\sigma_8=0.8159$, $n_s = 0.9667$).
The mass resolution in the high-resolution Lagrangian volume is
$m_\mathrm{DM}\approx 3.5\times10^4\,\Msun$ per dark matter particle.
Gas is represented at the same mass scale initially; the FIRE feedback
algorithms operate at this \textit{FIRE} resolution throughout the low-redshift
($z\lesssim 10$) evolution, but the Jeans-refinement module described in
Section~\ref{sec:jeans} further subdivides gas elements in dense star-forming
regions at $z\sim 21$.

The present study focuses exclusively on the redshift interval
$z\approx 21$ ($a\approx 0.045$, $t\approx 627\,\Myr$ after the Big Bang),
when the first protostellar collapse occurs inside a dense gas clump within
the proto-galactic halo.

\subsection{FIRE-3 physics}
\label{sec:fire3}

The FIRE-3 implementation of stellar physics and interstellar medium (ISM)
processes \citep{Hopkins2022} provides the sub-resolution source terms for
feedback and cooling at the FIRE resolution.

\paragraph{Radiative cooling and heating.}
Gas cools and heats through a comprehensive network spanning temperatures
$10 \lesssim T \lesssim 10^{10}\,\mathrm{K}$.  At high temperatures
($T > 10^4\,\mathrm{K}$) the cooling function includes free--free emission,
recombination, collisional ionisation and excitation of hydrogen and helium,
and metal-line cooling evaluated for 11 tracked metal species.  At
$T \lesssim 10^4\,\mathrm{K}$ and densities $n \gg 1\,\mathrm{cm}^{-3}$,
FIRE-3 adds fine-structure and rotational cooling from
C\,\textsc{ii}, O\,\textsc{i}, Si\,\textsc{ii}, and other low-ionisation
species, as well as molecular cooling from \ce{H2} and CO, an explicit
\ce{H2} formation and photodissociation network, and dust-gas thermal
coupling.  The metal cooling rates and chemical network are updated relative
to FIRE-2 to use revised low-temperature data, particularly important for
primordial and near-primordial gas at $z\sim 21$ where the metallicity is
essentially zero.

The extragalactic ionising background is included following
\citet{Faucher-Giguere2020} and is updated in FIRE-3 to match the most
recent observational constraints; at $z\sim21$ this background is negligible
compared with local protostellar and collisional processes.

\paragraph{Stellar evolution and feedback.}
FIRE-3 updates the stellar evolution tables relative to FIRE-2 to use
revised tracks from multiple stellar evolution codes, yielding mass-dependent
stellar luminosities, wind mass-loss rates, and supernova (SN) rates
consistent with current empirical constraints.  Feedback channels include:
(i) continuous radiation pressure from stellar spectra (UV through IR), coupled
self-consistently to the local gas opacity;
(ii) stellar winds from O/B stars and AGB stars, with rates and terminal
velocities interpolated from the updated evolutionary tracks;
(iii) core-collapse supernovae (CCSNe) following a delay-time distribution
tied to the FIRE-3 massive-star tracks, each depositing $E_\mathrm{SN} =
10^{51}\,\mathrm{erg}$ as thermal plus kinetic energy;
(iv) Type Ia supernovae with a delay-time distribution
$\mathrm{d}N/\mathrm{d}t \propto t^{-1}$ for $t > 40\,\mathrm{Myr}$.
Feedback is injected directly into the gas particles that overlap each star
particle's kernel, ensuring strict energy and momentum conservation.

At $z\sim 21$ the stellar population in the proto-galaxy is composed
exclusively of the first generation of metal-free (\textit{Population III})
stars, so the stellar feedback channels above operate in the primordial
limit: zero-metallicity cooling, enhanced ionising luminosities, and
potentially altered wind mass-loss rates.  The present work focuses on the
pre-supernova protostellar phase, so SN and AGB feedback play no role in
the epochs analysed.

\subsection{STARFORGE Jeans refinement}
\label{sec:jeans}

To follow gravitational collapse from galactic ($\sim\mathrm{kpc}$) to
protostellar ($\lesssim10\,\mathrm{AU}$) scales within the cosmological
simulation, we employ the \textsc{StarForge} module \citep{Grudic2021,
Grudic2022} as a ``Jeans-refinement'' overlay on the FIRE-3 gas.
\textsc{StarForge} extends \textsc{Gizmo} with:
(i) an adaptive particle-splitting algorithm that enforces the Jeans
criterion at all times;
(ii) sink particle formation and accretion;
and (iii) protostellar feedback.

\paragraph{Particle splitting and Jeans criterion.}
Gas particles are split into two equal-mass children whenever the
local Jeans mass $M_J$ falls below a fraction $f_J$ of the particle
mass $m_\mathrm{gas}$:
\begin{equation}
    M_J = \frac{\pi^{5/2}}{6}\frac{c_s^3}{\sqrt{G^3\,\rho}}
    < f_J\,m_\mathrm{gas},
\end{equation}
where $c_s = \sqrt{\gamma P/\rho}$ is the local sound speed, $\rho$ is the
gas density, and $f_J = 4$ is the adopted threshold
\citep{Grudic2021}.  Each split preserves the total mass, momentum,
energy, and magnetic flux of the parent particle, distributing children along
the local angular-momentum eigenvector to minimise numerical noise.  The
splitting is applied iteratively until the Jeans criterion is satisfied at
every refinement level.  Particle merging is performed in the inverse
situation (in expanding low-density regions) to prevent an unbounded growth
in particle count.

As a result, the effective mass resolution in the star-forming core reaches
$m_\mathrm{gas}\sim 10^{-3}$--$10^{-1}\,\Msun$ at the epochs studied here,
corresponding to a minimum Jeans-resolved cell length of order $\sim 30$
AU and a maximum resolvable density of
\begin{equation}
    \rho_J \approx 3\times10^{-14}\,\mathrm{g\,cm^{-3}}
    \left(\frac{m_\mathrm{gas}}{10^{-3}\,\Msun}\right)^{-2}
    \left(\frac{c_s}{0.2\,\mathrm{km\,s^{-1}}}\right)^{-6}.
    \label{eq:rhoJ}
\end{equation}
Once $\rho > \rho_J$, further collapse is captured by forming a sink
particle rather than continuing to split.

\paragraph{Sink particle formation.}
A sink (PartType5) is created when a gas particle simultaneously satisfies
all of the following criteria \citep{Grudic2021}:
\begin{enumerate}
    \item \textit{Density threshold:} $\rho > \rho_J$ (above the
          maximum Jeans-resolved density; Eq.~\ref{eq:rhoJ}).
    \item \textit{Convergent flow:} the local velocity divergence
          $\nabla\cdot\mathbf{v} < 0$ -- the gas is collapsing.
    \item \textit{Gravitational boundedness:} the kinetic plus thermal
          energy of the particle and its neighbours is less than their
          mutual gravitational potential energy, i.e.\ the virial
          parameter $\alpha_\mathrm{vir} < 1$.
    \item \textit{Local density maximum:} the particle is a density
          maximum within its kernel -- no denser neighbour exists within
          $r_\mathrm{sink}$.
\end{enumerate}
The sink accretion radius is $r_\mathrm{sink} = 18\,\mathrm{AU}\times
(m_\mathrm{gas}/10^{-3}\,\Msun)$.  Once formed, the sink accretes
neighbouring gas on a characteristic timescale
$t_\mathrm{acc} = 500\,\mathrm{yr}\times(m_\mathrm{gas}/10^{-3}\,\Msun)$
whenever a gas particle enters $r_\mathrm{sink}$ and is gravitationally
bound to the sink.  All gas accreted onto the sink is removed from the
simulation and its mass, momentum, and angular momentum are added to the
sink's conserved quantities.  Sinks interact gravitationally with all other
particles and with each other, allowing binary and multiple-star formation
to be followed self-consistently.

\paragraph{Protostellar evolution and feedback.}
Each sink is assigned a protostellar model following the
\citet{Offner2009} tracks, which track the internal luminosity, radius,
and surface temperature as a function of accreted mass.  Protostellar
feedback is injected as:
\begin{itemize}
    \item \textit{Protostellar jets:} bipolar collimated outflows launched
          at the Alfv\'en radius, carrying a fraction $f_w\sim 0.3$ of the
          accreted mass at velocities $v_j \sim v_\mathrm{Kep}(r_*)\sim
          30$--$100\,\mathrm{km\,s^{-1}}$.  New gas elements are
          spawned on-the-fly \citep{Grudic2021} at the desired
          velocity and mass-loading, avoiding the resolution loss inherent
          to purely Lagrangian approaches.
    \item \textit{Protostellar radiation:} the protostellar luminosity
          is injected into the radiation transport solver as a point source
          and propagated in five frequency bands (IR, optical, NUV, FUV,
          ionising), coupling dust and gas heating, radiation pressure,
          and photoionisation.
    \item \textit{Stellar winds} (once the star reaches the main sequence):
          mass-loss rate and terminal velocity from the FIRE-3 O-star tracks.
\end{itemize}
In the epochs analysed here ($t \lesssim 200\,\mathrm{kyr}$ after first
protostar formation) the stars are still in the protostellar/pre-main-sequence
phase, so feedback is dominated by jets and protostellar irradiation rather
than winds or supernovae.

\subsection{Spatial cutouts and analysis domain}
\label{sec:cutouts}

The full simulation snapshot at each output time contains
${\sim}\,8\times10^7$ gas particles spanning the entire zoom-in volume.
The vast majority of these are at the FIRE resolution
($m_\mathrm{gas}\sim 10^4\,\Msun$, kernel radii of order kpc)
and are irrelevant to the protostellar environment.  To enable efficient
analysis and visualisation, we extract spatial cutouts centered on the most
massive sink particle (or, before any sinks form, the densest Jeans-refined
gas particle) using a comoving aperture of $r_\mathrm{cut} = 0.005\,h^{-1}$
comoving kpc, corresponding to a physical radius of approximately
$325\,\mathrm{kAU} \approx 1.6\,\mathrm{pc}$ at $z\sim21$.
The cutout retains all PartType0 (gas), PartType3 (refinement tracers),
PartType4 (FIRE star particles), and PartType5 (STARFORGE sinks);
the $8.8\times10^7$ PartType1 dark matter particles are excluded to reduce
file size.

All analysis in this paper operates on these cutout snapshots.  The
Jeans-refined gas captured within each cutout accounts for essentially all
gas at densities relevant to protostellar disk formation
($\rho\gtrsim10^{-18}\,\mathrm{g\,cm^{-3}}$), and contains of order
$10^3$--$10^4$ particles per snapshot near the epoch of first sink
formation.

\subsection{Simulation parameters}
\label{sec:params}

Key numerical parameters are summarised in Table~\ref{tab:params}.

\begin{table}[h]
\centering
\caption{Key simulation parameters.}
\label{tab:params}
\begin{tabular}{lll}
\hline\hline
Parameter & Symbol & Value \\
\hline
\multicolumn{3}{l}{\textit{Cosmology}} \\
Matter density & $\Omega_\mathrm{m}$ & 0.3089 \\
Dark energy density & $\Omega_\Lambda$ & 0.6911 \\
Baryon density & $\Omega_\mathrm{b}$ & 0.0486 \\
Hubble constant & $H_0$ & $67.74\,\mathrm{km\,s^{-1}\,Mpc^{-1}}$ \\
\hline
\multicolumn{3}{l}{\textit{FIRE-3 resolution (initial)}} \\
DM particle mass & $m_\mathrm{DM}$ & $\sim 3.5\times10^4\,\Msun$ \\
Gas particle mass & $m_\mathrm{gas,FIRE}$ & $\sim 7\times10^3\,\Msun$ \\
\hline
\multicolumn{3}{l}{\textit{STARFORGE Jeans refinement}} \\
Jeans criterion & $f_J$ & 4 \\
Minimum resolved cell & $\ell_\mathrm{min}$ & $\sim 30\,\mathrm{AU}$ \\
Sink accretion radius & $r_\mathrm{sink}$ & $\sim 18\,\mathrm{AU}$ \\
Density threshold & $\rho_J$ & $\lesssim 10^{-14}\,\mathrm{g\,cm}^{-3}$ \\
\hline
\multicolumn{3}{l}{\textit{Analysis domain}} \\
Analysis epoch & $z$ & $\sim 21$ ($a\approx 0.045$) \\
Time of first sink & $t_1$ & $\approx 627.6\,\mathrm{Myr}$ \\
Cutout aperture (phys.) & $r_\mathrm{cut}$ & $\approx 325\,\mathrm{kAU}$ \\
Analysis box & $L_\mathrm{box}$ & $4\,\mathrm{kAU}$ (face-on maps) \\
Disk identification radius & $r_\mathrm{max}$ & $2\,\mathrm{kAU}$ \\
Number of snapshots & $N_\mathrm{snap}$ & 619 \\
\hline\hline
\end{tabular}
\end{table}
